\begin{abstract}
Large language model (LLM) agents are increasingly deployed over Model
Context Protocol (MCP) servers, yet the benchmarks used to evaluate them
score the final answer or a fixed ``ground-truth'' list of tools---both
of which are fragile once the underlying data is live and stateful. We
present \textbf{DynamicMCPBench}, a \emph{reusable framework} rather than
a fixed dataset. A practitioner can run it on their \emph{own} MCP servers
to test models on their own tasks, or let it collect servers automatically
to measure a model's \emph{general} ability to solve agentic tasks. Given
the servers and any set of models, it generates realistic goals, pursues
each one live to record a successful trajectory, distills that trajectory
into path-agnostic \emph{effect checkpoints}, and scores an agent on
whether it reproduces those effects---\emph{never} on the final
answer.\footnote{Code and data are released anonymously for review: code
at \url{https://anonymous.4open.science/r/DynamicMCPBench-4243/} and dataset at
\url{https://huggingface.co/datasets/anonsubmitter/DynamicMCPBench}.} To show what the framework reveals, we run it at scale:
24 models over 121 servers and 750 tasks spread evenly over 15 task
categories (50 each), where each category targets a distinct tool-use
challenge of the generated questions. Each task is scored by
pass\textasciicircum{}3---it counts as solved only if all three independent
attempts succeed. Even the strongest agents solve only about half of the
tasks, 31\% of tasks are solved by no model at all, and accuracy collapses
as the required tool chain grows longer (from 39\% on the shortest chains
to 13\% on the longest). A human validation study confirms the automatic
scoring is reliable (chance-corrected agreement of 0.76). DynamicMCPBench
thus turns benchmark construction into something practitioners can rerun on
their own servers and models, while exposing a consistent inability of
current agents to handle long, multi-step agentic tasks.
\end{abstract}
